\title{\centering \dsviv{}: \\ Towards Highly Efficient Million-Token Context Intelligence}

\author[*]{
\vspace{-0.8cm}
DeepSeek-AI
\\
\small
\vspace{-0.3cm}
\texttt{research@deepseek.com}
\vspace{-0.9cm}
}

\input{commands.tex}

\newcommand{\fix}{\marginpar{FIX}}
\newcommand{\new}{\marginpar{NEW}}

\newcommand{\zdaxie}[1]{\textcolor[rgb]{0.545,0,0.071}{#1}}

\begin{abstract}
\vspace{-0.5cm}

We present a preview version of \dsviv{} series, including two strong Mixture-of-Experts~(MoE) language models --- \dsvivp{} with \dsvivptp{} parameters (\dsvivpap{} activated) and \dsvivf{} with \dsvivftp{} parameters (\dsvivfap{} activated) --- both supporting a context length of one million tokens.
\dsviv{} series incorporate several key upgrades in architecture and optimization: 
(1) a hybrid attention architecture that combines \csafull{}~(\csa{}) and \hcafull{}~(\hca{}) to improve long-context efficiency; 
(2) \mhcfull{} (\mhc{}) that enhance conventional residual connections;
(3) and the Muon optimizer for faster convergence and greater training stability.
We pre-train both models on more than 32T diverse and high-quality tokens, followed by a comprehensive post-training pipeline that unlocks and further enhances their capabilities.
\dsvivpm{}, the maximum reasoning effort mode of \dsvivp{}, redefines the state-of-the-art for open models, outperforming its predecessors in core tasks.
Meanwhile, \dsviv{} series are highly efficient in long-context scenarios. 
In the one-million-token context setting, \dsvivp{} requires only 27\% of single-token inference FLOPs and 10\% of KV cache compared with \dsviiiII{}. 
This enables us to routinely support one-million-token contexts, thereby making long-horizon tasks and further test-time scaling more feasible.
The model checkpoints are available at \url{https://huggingface.co/collections/deepseek-ai/deepseek-v4}. 

\end{abstract}

\begin{document}
\begin{CJK*}{UTF8}{gbsn}

\maketitle

\begin{figure}[h]
\centering
\includegraphics[width=1.0\textwidth]{figures/dsv4_performance.pdf}
\caption{
    \textbf{Left}: benchmark performance of \dsvivpm{} and its counterparts.
    \textbf{Right}: inference FLOPs and KV cache size of \dsviv{} series and \dsviiiII. 
}
\label{fig:dsv4_performance_and_cost}
\end{figure}

\newpage

\begin{spacing}{0.9}
\tableofcontents
\end{spacing}

\newpage

\section{Introduction}

% background
The emergence of reasoning models~\citep{o1,dsr1} has established a new paradigm of test-time scaling, driving substantial performance gains for Large Language Models~(LLMs). 
However, this scaling paradigm is fundamentally constrained by the quadratic computational complexity of the vanilla attention mechanism~\citep{transformer}, which creates a prohibitive bottleneck for ultra-long contexts and reasoning processes. 
Concurrently, the emergence of long-horizon scenarios and tasks --- from complex agentic workflows to massive cross-document analysis --- has also made efficient support for ultra-long contexts critical for future progress. 
While recent open-source efforts~\citep{dsv3,qwen3,minimax_m2,kimi_k2} have advanced general capabilities, this core architectural inefficiency in handling ultra-long sequences remains a key impediment, limiting further gains from test-time scaling and hindering further exploration into long-horizon scenarios and tasks.

% brief intro to dsv4
In order to break the efficiency barrier in ultra-long contexts, we develop the \dsviv{} series, including the preview versions of \dsvivp{} with \dsvivptp{} parameters (\dsvivpap{} activated) and \dsvivf{} with \dsvivftp{} parameters (\dsvivfap{} activated). 
Through architectural innovations, \dsviv{} series achieve a dramatic leap in computational efficiency for processing ultra-long sequences. 
This breakthrough enables efficient support for a context length of one million tokens, ushering in a new era of million-length contexts for next-generation LLMs.
We believe our capability to efficiently handle ultra-long sequences unlocks the next frontier of test-time scaling, paves the way for deeper research into long-horizon tasks, and establishes a necessary foundation for exploring future paradigms like online learning.

% brief intro to arch
Compared with the \dsviii{} architecture~\citep{dsv3}, \dsviv{} series retain the DeepSeekMoE framework~\citep{deepseekmoe} and Multi-Token Prediction (MTP) strategy, while introducing several key innovations in architecture and optimization.  
To enhance long-context efficiency, we design a hybrid attention mechanism combining \csafull{}~(\csa{}) and \hcafull{}~(\hca{}). 
\csa{} compresses the KV caches along the sequence dimension and then performs DeepSeek Sparse Attention (DSA)~\citep{dsv32}, whereas \hca{} applies more aggressive compression to the KV caches but keeps dense attention.
To strengthen modeling capability, we incorporate \mhcfull{} (\mhc{})~\citep{mhc} that upgrade conventional residual connections.  
Additionally, we introduce the Muon~\citep{muon,muon_kimi} optimizer to the training of \dsviv{} series, leading to faster convergence and improved training stability.

% brief intro to infra
To enable efficient training and inference for \dsviv{} series as well as productive development, we introduce several infrastructure optimizations.
First, we design and implement a single fused kernel for MoE modules that fully overlaps computation, communication, and memory access. 
Second, we employ TileLang~\citep{wangtilelang}, a Domain-Specific Language (DSL) to balance development productivity and runtime efficiency. 
Third, we provide efficient batch-invariant and deterministic kernel libraries to ensure bitwise reproducibility across training and inference. 
Fourth, for the training framework, we extend the autograd framework with tensor-level checkpointing for fine-grained recomputation control; and we enhance training efficiency with a hybrid ZeRO strategy for the Muon optimizer, cost-effective \mhc{} implementations via recomputation and fused kernels, and two-stage contextual parallelism to manage compressed attention.
Fifth, for the inference framework, we design a heterogeneous KV cache structure with on-disk storage strategies to enable efficient shared-prefix reuse.
In addition, during the post-training stage, we incorporate FP4 quantization-aware training for MoE expert weights and the indexer QK path to reduce memory and computation. 

% efficiency and cost
By employing hybrid \csa{} and \hca{}, along with precision optimizations on computation and storage, \dsviv{} series achieve significantly lower inference FLOPs and a substantially reduced KV cache size compared with \dsviiiII{}, especially in long-context settings. 
The right part of Figure~\ref{fig:dsv4_performance_and_cost} demonstrates the estimated single-token inference FLOPs and accumulated KV cache size of \dsviiiII{} and \dsviv{} series.
In the scenario of 1M-token context, even \dsvivp{}, which has a larger number of activated parameters, attains only 27\% of the single-token FLOPs (measured in equivalent FP8 FLOPs) and 10\% of the KV cache size relative to \dsviiiII{}. 
Furthermore, \dsvivf{}, with its smaller number of activated parameters, pushes efficiency even further: in the 1M-token context setting, it achieves only 10\% of the single-token FLOPs and 7\% of the KV cache size compared with \dsviiiII{}. 
Additionally, for \dsviv{} series, the routed expert parameters utilize FP4 precision. 
While the peak FLOPs for FP4 $\times$ FP8 operations are currently the same as FP8 $\times$ FP8 on existing hardware, they can theoretically be implemented to be $1/3$ more efficient on future hardware, which will further enhance the efficiency of \dsviv{} series.

% brief intro to pre-training
During pre-training, we train \dsvivf{} on \dsvivftoken{} tokens and \dsvivp{} on \dsvivptoken{} tokens, respectively. 
After pre-training, these two models can natively and efficiently support 1M-length contexts. 
In our internal evaluations, \dsvivf{}-Base already surpasses \dsviiiII{}-Base across a majority of benchmarks with its more parameter-efficient design. 
\dsvivp{}-Base further extends this advantage to set a new performance standard among DeepSeek foundation models, achieving comprehensive superiority across reasoning, coding, long-context, and world knowledge tasks.

% brief intro to post-training
The post-training pipeline of \dsviv{} series features a two-stage paradigm: the independent cultivation of domain-specific experts, followed by unified model consolidation via on-policy distillation~\citep{lu2025onpolicydistillation,minillm}. 
Initially, for each target domain --- such as mathematics, coding, agent, and instruction following --- a separate expert model is trained independently. 
The base model first undergoes Supervised Fine-Tuning (SFT) on high-quality, domain-specific data to establish foundational capabilities. 
Subsequently, Reinforcement Learning (RL) is applied using Group Relative Policy Optimization (GRPO)~\citep{dsr1}, which further optimizes the model for domain-aligned behaviors guided by reward models tailored to specific success criteria. 
This phase yields a diverse set of specialized experts, each excelling in its respective field. 
Finally, to integrate these distinct proficiencies, a single unified model is trained through on-policy distillation, wherein the unified model acts as the student learning to optimize the reverse KL loss with teacher models.

\noindent
\textbf{Summary of Core Evaluation Results}
\begin{itemize}[topsep=0pt]
    \item \textbf{Knowledge}:
    In assessments of broad world knowledge, \dsvivpm{}, the maximum reasoning effort mode of \dsvivp{}, significantly outperforms leading open-source models on the SimpleQA~\citep{simpleqa} and Chinese-SimpleQA~\citep{csimpleqa} benchmarks. 
    Regarding educational knowledge --- evaluated via MMLU-Pro~\citep{mmlu_pro}, HLE~\citep{hle}, and GPQA~\citep{gpqa} --- \dsvivpm{} shows a marginal lead over its open-source counterparts. 
    \dsvivpm{} has significantly closed the gap with the leading proprietary model, Gemini-3.1-Pro, despite still trailing it in these knowledge-based evaluations.
    
    \item \textbf{Reasoning}: 
    Through the expansion of reasoning tokens, \dsvivpm{} demonstrates superior performance relative to GPT-5.2 and Gemini-3.0-Pro on standard reasoning benchmarks. 
    Nevertheless, its performance falls marginally short of GPT-5.4 and Gemini-3.1-Pro, suggesting a developmental trajectory that trails state-of-the-art frontier models by approximately 3 to 6 months.
    Furthermore, \dsvivfm{} achieves comparable performance to GPT-5.2 and Gemini-3.0-Pro, establishing itself as a highly cost-effective architecture for complex reasoning tasks.

    \item \textbf{Agent}: 
    On public benchmarks, \dsvivpm{} is on par with leading open-source models, such as Kimi-K2.6 and GLM-5.1, but slightly worse than frontier closed models. 
    In our internal evaluation, \dsvivpm{} outperforms Claude Sonnet 4.5 and approaches the level of Opus 4.5. 

    \item \textbf{Long-Context}: \dsvivpm{} delivers strong results on synthetic and real use cases with a 1-million-token context window, surpassing even Gemini-3.1-Pro on academic benchmarks.

    \item \textbf{\dsvivp{} v.s. \dsvivf{}}: 
    \dsvivfm{} exhibits lower performance in knowledge evaluations due to its smaller parameter scale. 
    However, it achieves comparable results on reasoning tasks when allocated a larger thinking budget. 
    In agent evaluations, while \dsvivfm{} matches the performance of \dsvivpm{} on several benchmarks, it still trails its larger counterpart on more complex, high-difficulty tasks.

\end{itemize}

\begin{figure}[t]
\centering
\includegraphics[width=0.6\linewidth]{figures/basic_arch.pdf}
\caption{
    Overall architecture of \dsviv{} series. 
    We use hybrid \csa{} (\csafull{}) and \hca{} (\hcafull{}) for attention layers, DeepSeekMoE for feed-forward layers, and strengthen conventional residual connections with \mhc{}.
}
\label{fig:dsv4_arch}
\end{figure}

\section{Architecture}
\label{sec:arch}

Overall, \dsviv{} series retain the Transformer~\citep{transformer} architecture and Multi-Token Prediction (MTP) modules~\citep{meta_mtp,dsv3}, while introducing several key upgrades over \dsviii{}:
(1) firstly, we introduce the \mhcfull{} (\mhc{})~\citep{mhc} to strengthen conventional residual connections; 
(2) secondly, we design a hybrid attention architecture, which greatly improves long-context efficiency through \csafull{} and \hcafull{}. 
(3) thirdly, we employ Muon~\citep{muon,muon_kimi} as the optimizer.
For the Mixture-of-Experts~(MoE) components, we still adopt the \dsmoe{}~\citep{deepseekmoe} architecture, with only minor adjustments from \dsviii{}.
The Multi-Token Prediction (MTP)~\citep{ms_mtp,meta_mtp,eagle,dsv3} configuration remains identical to that of \dsviii{}. 
All other unspecified details follow the settings established in \dsviii{}~\citep{dsv3}. 
Figure~\ref{fig:dsv4_arch} illustrates the overall architecture of \dsviv{}, and the details are described below.

\subsection{Designs Inherited from \dsviii{}}

\paragraph{Mixture-of-Experts.}

As previous DeepSeek-series models~\citep{dsvii,dsv3}, \dsviv{} series also adopt the \dsmoe{} paradigm~\citep{deepseekmoe} for Feed-Forward Networks~(FFNs), which sets fine-grained routed experts and shared experts.
Different from \dsviii{}, we change the activation function that computes the affinity scores from $\operatorname{Sigmoid}(\cdot)$ into $\operatorname{Sqrt}(\operatorname{Softplus}(\cdot))$.
For load balancing, we also employ the auxiliary-loss-free strategy~\citep{noaux_tc,dsv3}, augmented by a slight sequence-wise balance loss that prevents extreme imbalance within individual sequences.
For \dsviv{}, we remove the constraint on the number of routing target nodes, and carefully redesign the parallelism strategy to maintain training efficiency.
Furthermore, compared with \dsviii{}, we replace the dense FFN layers in the initial several Transformer blocks with MoE layers that employ Hash routing~\citep{hash_layer}. 
The Hash routing strategy determines the target experts of each token according to a predefined hash function with regard to the input token ID.

\paragraph{Multi-Token Prediction.}

As \dsviii{}, \dsviv{} series also set MTP modules and objectives. 
Given that the MTP strategy has been validated in \dsviii{}, we adopt the same strategy for \dsviv{} series without modification.

\subsection{\mhcfull{}}

As shown in Figure~\ref{fig:dsv4_arch}, \dsviv{} series incorporate \mhcfull{} (\mhc{})~\citep{mhc} to strengthen the conventional residual connections between adjacent Transformer blocks. 
Compared with naive Hyper-Connections (HC)~\citep{hc}, the core idea of \mhc{} is to constrain the residual mapping onto a specific manifold, and thus enhance the stability of signal propagation across layers while preserving model expressivity.
This subsection briefly introduces the standard HC and describes how we design \mhc{} for stable training. 

\paragraph{Standard Hyper-Connections.}
The standard HC expands the width of the residual stream by a factor of $n_{\text{hc}}$. 
Specifically, the shape of the residual stream is expanded from $\mathbb{R}^{d}$ to $\mathbb{R}^{n_{\text{hc}} \times d}$, where $d$ is the hidden size of the actual layer input. 
Let $X_l = [\mathbf{x}_{l,1}; \ldots; \mathbf{x}_{l,n_{\text{hc}}}]^T \in \mathbb{R}^{n_{\text{hc}} \times d}$ be the residual state before the $l$-th layer. 
HC introduces three linear mappings: an input mapping $A_{l} \in \mathbb{R}^{1 \times n_{\text{hc}}}$, a residual transformation $B_{l} \in \mathbb{R}^{n_{\text{hc}} \times n_{\text{hc}}}$, and an output mapping $C_{l} \in \mathbb{R}^{n_{\text{hc}} \times 1}$. 
The update of the residual state is then formulated as:
\begin{equation}
    X_{l+1} = B_{l} X_l + C_{l} \mathcal{F}_{l}(A_{l} X_l),
\end{equation}
where $\mathcal{F}_{l}$ denotes the $l$-th layer (e.g., an MoE layer), whose input and output shapes are both $\mathbb{R}^{d}$.
Note that the actual layer input $A_{l} X_l \in \mathbb{R}^{d}$ is also $d$-dimensional, so the expanded residual width does not influence the design of the inner layers. 
HC decouples the residual width from the actual hidden size, offering a complementary scaling axis with minimal computational overhead, as $n_{\text{hc}}$ is typically much smaller than the hidden size $d$.
However, even though HC has demonstrated potential in improving model performance, we find that the training will frequently exhibit numerical instability when stacking multiple layers, which hinders the scaling of HC.

\paragraph{Manifold-Constrained Residual Mapping.}
The core innovation of \mhc{} is to constrain the residual mapping matrix $B_l$ to the manifold of doubly stochastic matrices (the Birkhoff polytope) $\mathcal{M}$, and thus enhance the stability of signal propagation across layers:
\begin{equation}
    B_l \in \mathcal{M} \coloneq \{ M \in \mathbb{R}^{n \times n} \mid M\mathbf{1}_n = \mathbf{1}_n, \; \mathbf{1}_n^T M = \mathbf{1}_n^T, \; M \geq 0 \}.
\end{equation}
This constraint ensures that the spectral norm of the mapping matrix $\|B_l\|_2$ is bounded by 1, so the residual transformation is non-expansive, which increases the numerical stability during both the forward pass and backpropagation.
Besides, the set $\mathcal{M}$ is closed under multiplication, which guarantees stability in the scenarios of deep stacks of \mhc{}. 
In addition, the input transformation $A_l$ and output transformation $C_l$ are also constrained to be non-negative and bounded via a Sigmoid function to avoid the risk of signal cancellation.

\paragraph{Dynamic Parameterization.}
The parameters of three linear mappings are dynamically generated, which are decomposed into a dynamic (input-dependent) component and a static (input-independent) component. 
Given the input $X_l \in \mathbb{R}^{n_{\text{hc}} \times d}$, it is first flattened and normalized: $\hat{X}_l = \operatorname{RMSNorm}(\operatorname{vec}(X_l)) \in \mathbb{R}^{1 \times n_{\text{hc}}d}$. 
Then, we follow the conventional HC to generate the unconstrained raw parameters $\tilde{A}_l \in \mathbb{R}^{1 \times n_{\text{hc}}}$, $\tilde{B}_l \in \mathbb{R}^{n_{\text{hc}} \times n_{\text{hc}}}$, and $\tilde{C}_l \in \mathbb{R}^{n_{\text{hc}} \times 1}$:
\begin{align}
    \tilde{A}_l &= \alpha_l^\mathrm{pre} \cdot (\hat{X}_l W^\mathrm{pre}_l) + S_l^\mathrm{pre}, \\
    \tilde{B}_l &= \alpha_l^\mathrm{res} \cdot \operatorname{Mat}(\hat{X}_l W^\mathrm{res}_l) + S_l^\mathrm{res}, \\
    \tilde{C}_l &= \alpha_l^\mathrm{post} \cdot (\hat{X}_l W^\mathrm{post}_l)^T + S_l^\mathrm{post},
\end{align}
where $W^\mathrm{pre}_l, W^\mathrm{post}_l \in \mathbb{R}^{n_{\text{hc}}d \times n_{\text{hc}}}$ and $W^\mathrm{res}_l \in \mathbb{R}^{n_{\text{hc}}d \times n_{\text{hc}}^2}$ are learnable parameters for generating the dynamic components;
$\operatorname{Mat}(\cdot)$ reshapes a vector of size $1 \times n_{\text{hc}}^2$ into a matrix of size $n_{\text{hc}} \times n_{\text{hc}}$; 
$S_l^\mathrm{pre} \in \mathbb{R}^{1 \times n_{\text{hc}}}$, $S_l^\mathrm{post} \in \mathbb{R}^{n_{\text{hc}}\times 1}$, and $S_l^\mathrm{res} \in \mathbb{R}^{n_{\text{hc}} \times n_{\text{hc}}}$ are learnable static biases; 
and $\alpha_l^\mathrm{pre}$, $\alpha_l^\mathrm{res}$, $\alpha_l^\mathrm{post} \in \mathbb{R}$ are learnable gating factors initialized to small values. 

\paragraph{Applying Parameter Constraints.}
After obtaining the unconstrained raw parameters $\tilde{A}_l, \tilde{B}_l, \tilde{C}_l$, we then apply constraints described earlier to them to enhance the numerical stability. 
To be specific, for the input and output mappings, we employ a Sigmoid function $\sigma(\cdot)$ to ensure their non-negativity and boundedness:
\begin{align}
    A_l &= \sigma(\tilde{A}_l), \\
    C_l &= 2\sigma(\tilde{C}_l).
\end{align}
As for the residual mapping $\tilde{B}_l$, we project it onto the manifold of doubly stochastic matrices $\mathcal{M}$. 
This is achieved by the Sinkhorn-Knopp algorithm, which first applies an exponential function to $\tilde{B}_l$ to ensure positivity, getting $M^{(0)} = \exp(\tilde{B}_l)$, and then iteratively performs column and row normalization:
\begin{equation}
    M^{(t)} = \mathcal{T}_r(\mathcal{T}_c(M^{(t-1)})),
\end{equation}
where $\mathcal{T}_r$ and $\mathcal{T}_c$ denote row and column normalization, respectively. 
This iteration converges to a constrained doubly stochastic matrix $B_l = M^{(t_{\text{max}})}$. 
We choose $t_{\text{max}} = 20$ as a practical value.

\subsection{Hybrid Attention with \csa{} and \hca{}}

As the context length reaches extreme scales, the attention mechanism emerges as the dominant computational bottleneck in a model.
For \dsviv{}, we design two efficient attention architectures --- \csafull{}~(\csa{}) and \hcafull{}~(\hca{}) --- and employ their interleaved hybrid configuration, which substantially reduces the computational cost of attention in long-text scenarios. 
\csa{} integrates both compression and sparse attention strategies: it first compresses the Key-Value~(KV) cache of every $m$ tokens into one entry, and then applies \dsafull{} (\dsa{})~\citep{dsv32} where each query token attends to only $k$ compressed KV entries.
\hca{} aims for extreme compression by consolidating the KV cache of every $m^{\prime}$ ($\gg m$) tokens into a single entry. 
The hybrid architecture of \csa{} and \hca{} remarkably improves the long-context efficiency of \dsviv{} series, making one-million-token context feasible in practice. 
This subsection describes the core techniques of our hybrid attention architecture, and we also provide an open-source implementation\footnote{\url{https://huggingface.co/deepseek-ai/DeepSeek-V4-Pro/tree/main/inference}} to specify more details unambiguously. 

\begin{figure}[t]
\centering
\includegraphics[width=0.99\linewidth]{figures/CSA.pdf}
\caption{
    Core architectures of \csa{}. 
    It compresses the number of KV entries to $\frac{1}{m}$ times, and then applies \dsafull{} for further acceleration. 
    Additionally, a small set of sliding window KV entries is combined with the selected compressed KV entries to enhance local fine-grained dependencies. 
}
\label{fig:csa}
\end{figure}

\subsubsection{\csafull{}}

The core architecture of \csa{} is illustrated in Figure~\ref{fig:csa}, which first compresses the KV cache of each $m$ tokens into one entry, and then applies \dsafull{} for further acceleration. 

\paragraph{Compressed Key-Value Entries.}
Let $H \in \mathbb{R}^{n \times d}$ be a sequence of input hidden states, where $n$ is the sequence length and $d$ is the hidden size. 
\csa{} first computes two series of KV entries $C^{a}, C^{b} \in \mathbb{R}^{n \times c}$ and their corresponding compression weights $Z^{a}, Z^{b} \in \mathbb{R}^{n \times c}$, where $c$ is the head dimension:
\begin{align}
    C^{a} &= H \cdot W^{aKV}, \quad C^{b} = H \cdot W^{bKV}, \\
    Z^{a} &= H \cdot W^{aZ}, \quad~~ Z^{b} = H \cdot W^{bZ},
\end{align}
where $W^{aKV}, W^{bKV}, W^{aZ}, W^{bZ} \in \mathbb{R}^{d \times c}$ are trainable parameters.
Next, each $m$ KV entries in $C^{a}$ and $C^{b}$ will be compressed into one entry according to their compression weights and learnable positional biases $B^a,B^b \in \mathbb{R}^{m \times c}$, producing $C^{\text{Comp}} \in \mathbb{R}^{\frac{n}{m} \times c}$. 
Each compressed entry $C^{\text{Comp}}_{i} \in \mathbb{R}^{c}$ is computed by
\begin{align}
     [S^a_{mi:m(i+1)-1};S^b_{m(i-1):mi-1}] &=  \operatorname{Softmax}_{\text{row}}([Z^{a}_{mi:m(i+1)-1} + B^a;Z^{b}_{m(i-1):mi-1} + B^b]), \\
    C^{\text{Comp}}_{i} &= \sum_{j=mi}^{m(i+1)-1} S^a_j \odot C^{a}_{j} + \sum_{j=m(i-1)}^{mi-1} S^b_j \odot C^{b}_{j},
\end{align}
where $\odot$ denotes the Hadamard product; 
$\operatorname{Softmax}_{\text{row}}(\cdot)$ denotes the softmax operation along the row dimension, which performs normalization across the total of $2m$ elements from both $Z^a$ and $Z^b$. 
When $i=0$, $Z^b_{m(i-1):mi-1}$ is padded with negative infinity and $C^b_{m(i-1):mi-1}$ is padded with zeros.
Note that each $C^{\text{Comp}}_{i}$ is derived from $2m$ KV entries, but the indexes of $C^b$ used for $C^{\text{Comp}}_{i}$ and the indexes of $C^a$ used for $C^{\text{Comp}}_{i-1}$ are overlapped.
Therefore, \csa{} in fact compresses the sequence length to $\frac{1}{m}$ times. 

\paragraph{Lightning Indexer for Sparse Selection.}
After obtaining the compressed KV entries $C^{\text{Comp}}$, \csa{} applies the \dsa{} strategy to select top-k compressed KV entries for core attention. 
First, \csa{} performs the same compression operation used for $C^{\text{Comp}}$ to get compressed indexer keys $K^{\text{IComp}} \in \mathbb{R}^{\frac{n}{m} \times c^I}$, where $c^I$ is the indexer head dimension. 
Then, for a query token $t$, we produce the indexer queries $\{\mathbf{q}_{t, 1}^{I};\mathbf{q}_{t, 2}^{I};...;\mathbf{q}_{t, n_{h}^{I}}^{I}\}$ in a low-rank manner:
\begin{align}
    \mathbf{c}_{t}^{Q} &= \mathbf{h}_{t} \cdot W^{DQ}, \\
    [\mathbf{q}_{t, 1}^{I};\mathbf{q}_{t, 2}^{I};...;\mathbf{q}_{t, n_{h}^{I}}^{I}] = \mathbf{q}_{t}^{I} &= \mathbf{c}_{t}^{Q} \cdot W^{IUQ},
\end{align}
where $\mathbf{h}_{t} \in \mathbb{R}^{d}$ is the input hidden state of the query token $t$; $\mathbf{c}_{t}^{Q} \in \mathbb{R}^{d_c}$ is the compressed latent vector for queries; $d_c$ denotes the query compression dimension; $n_{h}^{I}$ denotes the number of indexer query heads; $W^{DQ} \in \mathbb{R}^{d \times d_c}$ and $W^{IUQ} \in \mathbb{R}^{d_c \times c^I n_h^I}$ are the down-projection and up-projection matrices for indexer queries, respectively.
Next, the index score $I_{t, s} \in \mathbb{R}$ between the query token $t$ and a preceding compressed block $s$ ($s$ < $\operatorname{Floor}(\frac{t}{m})$) is computed by
\begin{align}
    [w_{t, 1}^I; w_{t, 2}^I; ...; w_{t, n_{h}^{I}}^I] = \mathbf{w}_t^I & = \mathbf{h}_{t} \cdot W^w, \\
    I_{t, s} & = \sum_{h=1}^{n_h^I} w_{t, h}^I \cdot \text{ReLU}\left(\mathbf{q}^{I}_{t, h} \cdot K^{\text{IComp}}_{s}\right),
\end{align}
where $W^w \in \mathbb{R}^{d \times n_{h}^{I}}$ is a learnable matrix; $w_{t, h}^I \in \mathbb{R}$ is the weight of the $h$-th indexer head.
For a query token $t$, given its index scores $I_{t, :}$, we employ a top-k selector to selectively retain a subset of compressed KV entries $\mathcal{C}^{\text{SprsComp}}_t$ for subsequent core attention:
\begin{equation}
    \mathcal{C}^{\text{SprsComp}}_t = \left\{ C^{\text{Comp}}_{s} ~\Big|~ I_{t, s} \in \operatorname{Top-k} (I_{t, :}) \right\}.
\end{equation}

\paragraph{Shared Key-Value MQA.}
After selecting the sparse KV entries, \csa{} then performs core attention in a Multi-Query Attention~(MQA)~\citep{mqa} manner, where each compressed KV entry in $\mathcal{C}^{\text{SprsComp}}_t$ serves as both attention key and value. 
To be specific, for a query token $t$, we first produce attention queries $\{\mathbf{q}_{t, 1};\mathbf{q}_{t, 2};...;\mathbf{q}_{t, n_{h}}\}$ from the compressed latent vector $\mathbf{c}_{t}^{Q}$:
\begin{equation}
    [\mathbf{q}_{t, 1};\mathbf{q}_{t, 2};...;\mathbf{q}_{t, n_{h}}] = \mathbf{q}_{t} = \mathbf{c}_{t}^{Q} \cdot W^{UQ},
\end{equation}
where $n_{h}$ denotes the number of query heads; $W^{UQ} \in \mathbb{R}^{d_c \times c n_h}$ is the up-projection matrices for queries.
Note that the latent query vector $\mathbf{c}_{t}^{Q}$ is shared with that used for the indexer queries. 
Next, we perform MQA on $\{\mathbf{q}_{t, i}\}$ and $\mathcal{C}^{\text{SprsComp}}_t$: 
\begin{equation}
    \mathbf{o}_{t,i} = \operatorname{CoreAttn}\left( \texttt{query=}\mathbf{q}_{t,i}, \texttt{key=}\mathcal{C}^{\text{SprsComp}}_t, \texttt{value=}\mathcal{C}^{\text{SprsComp}}_t \right),
\end{equation}
where $\mathbf{o}_{t,i} \in \mathbb{R}^{c}$ is the core attention output of the $i$-th head at the $t$-th token; $\operatorname{CoreAttn}(\cdot)$ denotes the core attention operation. 

\paragraph{Grouped Output Projection.}
In the configuration of \dsviv{}, $c n_h$ is quite large.
Therefore, directly projecting the outputs of the core attention operation $[\mathbf{o}_{t,1}; \mathbf{o}_{t,2}; ...; \mathbf{o}_{t,n_h}]=\mathbf{o}_{t} \in \mathbb{R}^{c n_h}$ to a $d$-dimensional hidden state will impose a substantial computational burden. 
To mitigate this cost, we design a grouped output projection strategy.
To be specific, we first split $n_h$ outputs into $g$ groups, and then for each group of output $\mathbf{o}^G_{t, i} \in \mathbb{R}^{c \frac{n_h}{g}}$, we project it to a $d_g$-dimensional intermediate output $\mathbf{o}^{G'}_{t, i} \in \mathbb{R}^{d_g}$, where $d_g < c \frac{n_h}{g}$. 
Finally, we project the intermediate output $[\mathbf{o}^{G'}_{t, 1}; \mathbf{o}^{G'}_{t, 2}; ...; \mathbf{o}^{G'}_{t, g}] \in \mathbb{R}^{d_gg}$ to the final attention output $\mathbf{\hat{o}}_{t} \in \mathbb{R}^{d}$. 

\begin{figure}[t]
\centering
\includegraphics[width=0.6\linewidth]{figures/HCA.pdf}
\caption{
    Core architectures of \hca{}.
    It performs heavier compression, where the KV entries of $m^{\prime}$ ($\gg m$) tokens will be consolidated into one. 
    Also, we additionally introduce a small set of sliding window KV entries to enhance local fine-grained dependencies. 
}
\label{fig:hca}
\end{figure}

\subsubsection{\hcafull{}}

The core architecture of \hca{} is illustrated in Figure~\ref{fig:hca}, which compresses the KV cache in a heavier manner, but does not employ sparse attention. 

\paragraph{Compressed Key-Value Entries.}
By and large, the compression strategy of \hca{} is similar to that of \csa{}, but employs a larger compression rate $m^{\prime} $ ($\gg m$) and does not perform overlapped compression. 
Let $H \in \mathbb{R}^{n \times d}$ be a sequence of input hidden states, \hca{} first computes the original KV entries $C \in \mathbb{R}^{n \times c}$ and their corresponding compression weights $Z \in \mathbb{R}^{n \times c}$:
\begin{align}
    C &= H \cdot W^{KV}, \\
    Z &= H \cdot W^{Z},
\end{align}
where $W^{KV}, W^{Z} \in \mathbb{R}^{d \times c}$ are trainable parameters.
Next, each $m^{\prime}$ KV entries in $C$ will be compressed into one according to the compression weights and learnable positional biases $B \in \mathbb{R}^{m^{\prime} \times c}$, producing $C^{\text{Comp}} \in \mathbb{R}^{\frac{n}{m^{\prime}} \times c}$. 
Each compressed entry $C^{\text{Comp}}_{i} \in \mathbb{R}^{c}$ is computed by
\begin{align}
     S_{m^{\prime}i:m^{\prime}(i+1)-1} &=  \operatorname{Softmax}_{\text{row}}(Z_{m^{\prime}i:m^{\prime}(i+1)-1} + B), \\
    C^{\text{Comp}}_{i} &= \sum_{j=m^{\prime}i}^{m^{\prime}(i+1)-1} S_j \odot C_{j}.
\end{align}
Through this compression operation, \hca{} compresses the sequence length to $\frac{1}{m^{\prime}}$ times. 

\paragraph{Shared Key-Value MQA and Grouped Output Projection.}
\hca{} also employs the shared KV MQA and grouped output projection strategies as \csa{} does. 
After the KV compression, for a query token $t$, \hca{} first produces attention queries $\{\mathbf{q}_{t, 1};\mathbf{q}_{t, 2};...;\mathbf{q}_{t, n_{h}}\}$ in a low-rank manner: 
\begin{align}
    \mathbf{c}_{t}^{Q} &= \mathbf{h}_{t} \cdot W^{DQ}, \\
    [\mathbf{q}_{t, 1};\mathbf{q}_{t, 2};...;\mathbf{q}_{t, n_{h}}] = \mathbf{q}_{t} &= \mathbf{c}_{t}^{Q} \cdot W^{UQ},
\end{align}
where $\mathbf{h}_{t} \in \mathbb{R}^{d}$ is the input hidden state of the query token $t$; $n_{h}$ denotes the number of query heads; $W^{DQ} \in \mathbb{R}^{d \times d_c}$ and $W^{UQ} \in \mathbb{R}^{d_c \times c n_h}$ are the down-projection and up-projection matrices for queries, respectively. 
Next, we perform MQA on $\{\mathbf{q}_{t, i}\}$ and $C^{\text{Comp}}$: 
\begin{equation}
    \mathbf{o}_{t,i} = \operatorname{CoreAttn}\left( \texttt{query=}\mathbf{q}_{t,i}, \texttt{key=}C^{\text{Comp}}, \texttt{value=}C^{\text{Comp}} \right),
\end{equation}
where $\mathbf{o}_{t,i} \in \mathbb{R}^{c}$ is the core attention output of the $i$-th head at the $t$-th token. 
Next, as \csa{} does, \hca{} splits $n_h$ outputs into $g$ groups, and for each group of output $\mathbf{o}^G_{t, i} \in \mathbb{R}^{c \frac{n_h}{g}}$, \hca{} projects it to a $d_g$-dimensional intermediate output $\mathbf{o}^{G'}_{t, i} \in \mathbb{R}^{d_g}$, where $d_g < c \frac{n_h}{g}$. 
Finally, \hca{} projects the intermediate output $[\mathbf{o}^{G'}_{t, 1}; \mathbf{o}^{G'}_{t, 2}; ...; \mathbf{o}^{G'}_{t, g}] \in \mathbb{R}^{d_gg}$ to the final attention output $\mathbf{\hat{o}}_{t} \in \mathbb{R}^{d}$. 

\subsubsection{Other Details}

In addition to the core architectures of \csa{} and \hca{} described above, our hybrid attention incorporates several other techniques. 
For writing clarity, we omit these additional techniques from the above introduction and will briefly describe them in this subsection.
Also, this subsection focuses only on the core ideas of them and may omit some tiny details for simplicity. 
We encourage the readers to refer to our open-source implementation for unambiguous details.

\paragraph{Query and Key-Value Entry Normalization.}
For both \csa{} and \hca{}, we perform an additional RMSNorm operation on each head of the queries and the only head of the compressed KV entries, just before the core attention operation. 
This normalization avoids exploding attention logits and may improve training stability. 

\paragraph{Partial Rotary Positional Embedding.}
For both \csa{} and \hca{}, we partially employ the Rotary Positional Embedding (RoPE)~\citep{su2024roformer} to the attention queries, KV entries, and the core attention outputs. 
To be specific, for each query vector and KV entry vector used in \csa{} and \hca{}, we apply RoPE to its last 64 dimensions. 
Since the KV entries serve as both attention keys and values, the naive core attention outputs $\{\mathbf{o}_{t,i}\}$ will carry absolute position embeddings, derived from the weighted sum of KV entries. 
As a countermeasure, we also apply RoPE with position $-i$ on the last 64 dimensions of each $\mathbf{o}_{t,i}$. 
In this way, the output of the core attention will also carry relative position embeddings --- the contribution of each KV entry to the core attention outputs will also be related to the distance between the query and the KV entry.

\paragraph{Additional Branch of Sliding Window Attention.}
In order to strictly preserve causality in \csa{} and \hca{}, each query attends to only preceding compressed KV blocks. 
Consequently, a query cannot access information from other tokens within its own compressed block.
Meanwhile, recent tokens usually possess greater relevance to the query token in language modeling.
For these reasons, we introduce a supplementary attention branch to both \csa{} and \hca{} in a sliding window manner, for better modeling of local dependencies.
To be specific, for each query token, we additionally produce $n_{\text{win}}$ uncompressed KV entries corresponding to the recent $n_{\text{win}}$ tokens. 
In the core attention of \csa{} and \hca{}, these KV entries in the sliding window will be used along with the compressed KV entries. 

\paragraph{Attention Sink.}
In the core attention of \csa{} and \hca{}, we employ the trick of attention sink~\citep{attn_sink,gpt_oss}. 
To be specific, we set a series of learnable sink logits $\{z^{\prime}_1, z^{\prime}_2, ..., z^{\prime}_{n_h}\}$. 
For the $h$-th attention head, $\operatorname{Exp}(z^{\prime}_h)$ will be added to the denominator of the attention score:
\begin{equation}
    s_{h, i, j} = \frac{\operatorname{Exp}(z_{h, i, j})}{\sum_k \operatorname{Exp}(z_{h, i, k}) + \operatorname{Exp}(z^{\prime}_h)},
\end{equation}
where $s_{h, i, j}, z_{h, i, j} \in \mathbb{R}$ denote the attention score and attention logit of the $h$-th attention head between the $i$-th query token and the $j$-th preceding token or compressed block.
This technique allows each query head to adjust its total attention scores to be not equal to 1, and even to be near 0.

\subsubsection{Efficiency Discussion}

Due to the employment of hybrid \csa{} and \hca{}, together with low-precision computation and storage, the attention module of \dsviv{} series achieves remarkable efficiency in both attention FLOPs and KV cache size, especially in long-context scenarios. 
First, we adopt a mixed storage format for KV entries: BF16 precision is used for the rotary positional embedding (RoPE) dimensions, while FP8 precision is applied to the remaining dimensions. 
This hybrid representation reduces the KV cache size by nearly half compared with pure BF16 storage. 
Second, attention computation within the lightning indexer is performed in FP4 precision, which accelerates the attention operation under extremely long contexts. 
Third, relative to \dsviiiII{}, a smaller attention top-k is chosen in \dsviv{} series, thereby improving model efficiency on short- and medium-length texts. 
Finally, and most importantly, compressed attention and hybrid attention techniques substantially reduce both the KV cache size and the computational FLOPs. 

Taking BF16 GQA8~\citep{ainslie2023gqa} with a head dimension of 128 as the baseline --- one of the common configurations of LLM attention --- the KV cache size of \dsviv{} series can be dramatically reduced to approximately $2\%$ times of that baseline in the 1M-context setting.
Moreover, even when compared with \dsviiiII{}~\citep{dsv32} --- already an efficient baseline --- \dsviv{} series still exhibits substantial advantages in efficiency. 
A comparison of their inference FLOPs and KV cache size is provided in the right part of Figure~\ref{fig:dsv4_performance_and_cost}.

\begin{algorithm}[t]
\caption{Muon Optimizer for \dsviv{}}
\label{alg:muon}
\begin{algorithmic}[1]
\Require{Learning rate $\eta$, momentum $\mu$, weight decay $\lambda$, update rescaling factor $\gamma$}
\For{each training step $t$}
    \For{each logically independent weight $W \in \mathbb{R}^{n \times m}$}
        \State $G_t = \nabla_{W} \mathcal{L}_t(W_{t-1})$ 
        \Comment{Compute gradients}
        \State $M_t = \mu M_{t-1} + G_t$ 
        \Comment{Accumulate momentum buffer}
        \State $O^{\prime}_t = \operatorname{HybridNewtonSchulz}(\mu M_t + G_t)$ 
        \Comment{Nesterov trick and hybrid Newton-Schulz}
        \State $O_t = O^{\prime}_t \cdot \sqrt{\max(n, m)} \cdot \gamma$ 
        \Comment{Rescale the update RMS}
        \State $W_t = W_{t-1} \cdot \left( 1 - \eta \lambda \right) - \eta O_t$ 
        \Comment{Perform weight decay and update}
    \EndFor
\EndFor
\end{algorithmic}
\end{algorithm}

\subsection{Muon Optimizer}

We employ the Muon~\citep{muon,muon_kimi} optimizer for the majority of modules in \dsviv{} series due to its faster convergence and improved training stability.
The full algorithm of our Muon optimization is summarized in Algorithm~\ref{alg:muon}.

\paragraph{Basic Configurations.}
We maintain the AdamW~\citep{adamW} optimizer for the embedding module, the prediction head module, the static biases and gating factors of \mhc{} modules, and the weights of all RMSNorm modules. 
All other modules are updated with Muon.
Following \citet{muon_kimi}, we also apply weight decay to Muon parameters, use the Nesterov~\citep{nesterov,muon} trick, and rescale the Root Mean Square (RMS) of the update matrix for reutilization of our AdamW hyper-parameters.
Different from them, we use hybrid Newton-Schulz iterations for orthogonalization.

\paragraph{Hybrid Newton-Schulz Iterations.}
For a given matrix $M$, let its Singular Value Decomposition (SVD) be $M = U \Sigma V^T$. 
The Newton-Schulz iterations aim to approximately orthogonalize $M$ to be $UV^T$.
Usually, $M$ will be first normalized as $M_0=M / ||M||_F$ to ensure its maximum singular value does not exceed 1. 
Then, each Newton-Schulz iteration performs the following operation:
\begin{equation}
    M_k = a M_{k-1} + b (M_{k-1} M_{k-1}^T) M_{k-1} + c (M_{k-1} M_{k-1}^T)^2 M_{k-1}.
\end{equation}
Our hybrid Newton-Schulz performs 10 iterations over two distinct stages. 
During the first 8 steps, we use coefficients $(a, b, c) = (3.4445, -4.7750, 2.0315)$ to drive rapid convergence, bringing the singular values close to 1. 
In the final 2 steps, we switch to coefficients $(a, b, c) = (2, -1.5, 0.5)$, which stabilize the singular values precisely at 1.

\paragraph{Avoiding Exploding Attention Logits.}
The attention architecture of \dsviv{} series allows us to directly apply RMSNorm on the attention queries and KV entries, which effectively prevents attention logits from exploding. 
Consequently, we do not employ the QK-Clip technique~\citep{muon_kimi} in our Muon optimizer. 
